\documentclass[aspectratio=169,xcolor=table]{beamer}
\usepackage{beamerthemesplit}
\usepackage{wrapfig}
\usetheme{SPbGU}
\usepackage{pdfpages}
\usepackage{amsmath}
\usepackage{cmap}
\usepackage[T2A]{fontenc}
\usepackage[utf8]{inputenc}
\usepackage[english]{babel}
\usepackage{indentfirst}
\usepackage{amsmath}
\usepackage{tikz}
\usepackage{multirow}
\usepackage[noend]{algpseudocode}
\usepackage{algorithm}
\usepackage{algorithmicx}
\usepackage{fancyvrb}
\usepackage{xcolor}
%usepackage{fancyvrb}
%\usepackage{minted}
%\usepackage{verbments}

\def\mathunderline#1#2{\color{#1}\underline{{\color{black}#2}}\color{black}}

\beamertemplatenavigationsymbolsempty

\title[CFPQ $\rightarrow$ ME]{Context-Free Path Querying via Matrix Equations}
\institute[SPbU]{
JetBrains Research, Programming Languages and Tools Lab  \\
Saint-Petersburg State University
}

\author[Yuliya Susanina]{Yuliya Susanina}

\date{14.06.2020}

\begin{document}
{
\begin{frame}[fragile]
  \begin{tabular}{p{2.0cm} p{7.5cm} p{1cm}}
   \begin{center}
      \includegraphics[height=1.5cm]{pictures/jetbrainsResearch.pdf}
    \end{center}
    &
    \begin{center}
      \includegraphics[height=1.5cm]{pictures/YC_logo.pdf}
    \end{center}
    &
    \begin{center}
      \includegraphics[height=1.5cm]{pictures/SPbGU_Logo.png}
    \end{center}
  \end{tabular}
  \titlepage
\end{frame}
}

\begin{frame}[fragile]
  \transwipe[direction=90]
  \frametitle{Context-Free Path Querying (CFPQ)}

\begin{itemize}
  \item Context-free grammar $G=(N, \Sigma, R)$ \\
  $\mathcal{L}(G_S) = \{\omega \mid S \Rightarrow_{G}^{*} \omega\}$, $S \in N$
  \item Directed graph $D = (V,E,\sigma)$, $\sigma \subseteq \Sigma$, $E\subseteq V\times \sigma \times V$ \\
  $m\lambda n$ --- path from $m$ to $n$ in $D$, $\lambda$ --- a unique word of this path
  \pause
  \only<2>{\item $\mathunderline{red} { R_A = \{(m, n) \mid m\lambda n \text{ --- path in } D, \lambda \in \mathcal{L}(G_A)\} }$}
  \pause
  \only<3>{\item $ R_A = \{(m, n) \mid m\lambda n \text{ --- path в } D, \lambda \in \mathcal{L}(G_A)\} $}
  \vspace{10pt}
  \item Application areas:
      \begin{itemize}
          \item Graph databases
          \item Bioinformatics
          \item Static code analisys
      \end{itemize}
\end{itemize}
\end{frame}

\begin{frame}[fragile]
  \transwipe[direction=90]
  \frametitle{Computational mathematics}
\begin{itemize}
  \item Both rich theoretical foundations and constantly improving implementations
  \item Parallel techniques and GPGPU
  \item Approximate computational methods
\end{itemize}
\pause
    \vspace{5pt}
    \color{red} {$$\Downarrow$$}
    \large\centering\color{red}{Acceleration of CFPQs processing}
   
\end{frame}


\begin{frame}[fragile]
  \transwipe[direction=90]
  \frametitle{Reduction from Solving Boolean Matrix Equations}
  
  $$S \rightarrow aSb \mid ab$$
  $$T_E \in \mathbb{M}^{|V| \times |V|}: (T_{E})_{ij} = 1 \iff (i,j) \in R_E \; \forall E \in (N \cup \Sigma)$$
  \pause
  $$\Downarrow$$
    \centering
    \begin{tabular}{cl}
    \parbox{0.07\linewidth}{
    $$\{ T_S^k \}:$$
    }
    & \begin{tabular}{l}
    $T_S^0 = \mathbf{0}$ \\
    $T_S^{k+1} = T_a T_S^k T_b + T_a T_b$
    \end{tabular} \\
    \end{tabular}

  $$T_S^{\infty} - \text{least solution } T_S = T_aT_ST_b + T_aT_b$$

\end{frame}

\begin{frame}[fragile]
  \transwipe[direction=90]
  \frametitle{Reduction from Solving Matrix Equations over $\mathbb{R}$}
  
    $$\Downarrow$$
    \centering
    \begin{tabular}{cl}
    \parbox{0.07\linewidth}{
    $\{ \mathcal{T}_S^k \}:$
    }
    & \begin{tabular}{l}
    $\mathcal{T}_S^0 = \mathbf{0}$ \\
    $\mathcal{T}_S^{k+1} = \epsilon(T_a \mathcal{T}_S^k T_b + T_a T_b)$
    \end{tabular} \\
    \end{tabular}

    \vspace{20pt}
    \(
    \left.
    \begin{array}{c}
    \mathcal{T}_S^{\infty} - \text{least solution } \mathcal{T}_S = \epsilon(T_a \mathcal{T}_S T_b + T_a T_b), \\
    \text{where } \epsilon \text{ such that } \mathcal{T}_{S}^{k} \leq \textbf{1} \quad \forall k
    \end{array}
    \right.
    \)
  
  \pause
  
  $$(\mathcal{T}_S^{k+1})_{ij} > 0 \iff (T_S^{k+1})_{ij} = 1$$
  $$ceil(\mathcal{T}_S^{\infty}) = T_S^{\infty}$$
\end{frame}

\begin{frame}[fragile]
  \transwipe[direction=90]
  \frametitle{Methods for Solving Equations}
 
\begin{itemize}
    \item Linear equations
    \begin{itemize}
        \item Sylvester equations $AXB + CXD = F$
        \item Linear systems $Ax=b$
    \end{itemize}
    \item Nonlinear equations
    \begin{itemize}
        \item Newton's method \\
                $$X = G(X) \Rightarrow F(X) = X - G(X) = \mathbf{0}$$
                \(
                \left.
                \begin{array}{l}
                X_{i+1} = X_i - (F'(X_i))^{-1}F(X_i)
                \end{array}
                \right.
                \iff
                \left\{
                \begin{array}{l}
                F'(X_i)H_i = - F(X_i) \\
                X_{i+1} = X_i + H_i
                \end{array}
                \right.
                \)
    \end{itemize}
\end{itemize}
\end{frame}

\begin{frame}[fragile]
  \transwipe[direction=90]
  \frametitle{First implementation}

\begin{itemize}
    \item \textit{SciPy}
    \begin{itemize}
        \item sSLV --- to solve as a sparse linear system
        \item dNWT --- to find a root of a function with Newton's method
    \end{itemize}
    \item Comparative analysis of matrix-based approach and equation-based approach (in ms)

        \begin{table}[h]
        \centering
        \begin{tabular}{ | c | c || c | c || c | c | c |}
        \hline
        Ontology    & |V| & dNWT & sSLV & dGPU & sCPU &  sGPU \\
        \hline
        \hline
        bio-meas    & 341 &  284 & 35   & 276  & 91  & 24\\
        people-pets & 337 &  73  & 49   & 144  & 38  & 6\\
        funding     & 778 &  502 & 184  & 1246 & 344 & 27\\
        wine        & 733 &  791 & 171  & 722  & 179 & 6\\
        pizza       & 671 &  334 & 161  & 943  & 256 & 23\\
        \hline
        \end{tabular}
        
        \end{table}

\end{itemize}
\end{frame}



\begin{frame}[fragile]
  \transwipe[direction=90]
  \frametitle{Results}
  
\begin{itemize}

  \item Equation-based approach for CFPQ was proposed
  \item The possibilities of using both accurate and approximate methods of computational mathematics was reviewed
  \item The evaluation on a set of conventional benchmarks showed that our approach is comparable with the matrix-based approach and applicable for real-world data processing
  
\end{itemize}

\end{frame}

\begin{frame}[fragile]
  \transwipe[direction=90]
  \frametitle{Future Work}

\begin{itemize}
    \item Employ high-performance solvers which utilize GPGPU and distributed computations
    \item Determine the subclasses of (system of) polynomial equations the solution of which can be reduced to CFPQ
    \item Try to construct a bidirectional reduction between CFPQ and these subclasses, thereby finding efficient solutions for both these problems
\end{itemize}

\end{frame}

\end{document}